\subsection{Adjoint}

In this section, K is one of the fields R or C, and E and F denote Hilbert spaces over K.

Let \(u\) be a densely defined operator from E to F. Let \(D\) denote the domain of \(u\) . For \(y \in \mathrm{F}\) , let \(\lambda_y\) the linear form on D given by \(\lambda_y(x) = \langle y \mid u(x) \rangle\) for all \(x \in \mathrm{D}\) . We denote \(\mathrm{D}^*\) the set of vectors \(y \in \mathrm{F}\) such that \(\lambda_y\) is continuous on D. It is a vector subspace of F. Let \(y \in \mathrm{D}^*\) ; since D is dense in E, the linear form \(\lambda_y\) extends uniquely to a continuous linear form on E, which we still denote by \(\lambda_y\) . According to EVT, V, p. 15, th. 3, there exists a unique element \(u^*(y)\) in E such that \(\lambda_y(x) = \langle u^*(y) \mid x \rangle\) for all \(x \in \mathrm{E}\) . The map \(y \mapsto u^*(y)\) is linear from \(\mathrm{D}^*\) in E.

DEFINITION 6. --- The partial operator from F to E with domain \(D^*\) defined by \(y \mapsto u^*(y)\) is called the adjoint of \(u\) . It is denoted \(u^*\) .

Remark. --- We therefore have \(y \in D^{*}\) if and only if there exists \(c \in R_{+}\) such that \(|\langle y | u(x) \rangle| \leqslant c \|x\|\) for all \(x \in D\) The element \(u^{*}(y)\) is then characterized by the relation

\[
\langle y \mid u (x) \rangle = \langle u ^ {*} (y) \mid x \rangle\tag{1}
\]

for all \(x \in \mathrm{D}\) . We then have \(|\langle y | u(x) \rangle| \leqslant \| u^{*}(y) \| \| x \|\) for all \(x \in \mathrm{D}\) .

In the case where \(u\) is a continuous linear map from E into F, its adjoint in the sense of the preceding definition coincides with the adjoint defined in EVT, V, p. 38, def. 1, since D* is equal to F in this case.

Let \(v \in \mathcal{P}(\mathrm{E};\mathrm{F})\) such that \(u \subset v\) . We then have \(v^{*} \subset u^{*}\) .

Let \(v \in \mathcal{P}(\mathrm{E};\mathrm{F})\) such that \(u + v\) has dense domain. The partial operator \(v\) then has dense domain and one has \(u^{*} + v^{*} \subset (u + v)^{*}\) . In general, there is no equality (exercise 9 of IV, p. 347). If \(v \in \mathcal{L}(\mathrm{E};\mathrm{F})\) , then \(u + v\) is densely defined and \((u + v)^{*} = u^{*} + v^{*}\) . This is the case, for example, if \(\mathrm{F} = \mathrm{E}\) and if \(v = \lambda 1_{\mathrm{E}}\) where \(\lambda \in \mathrm{K}\) .

Let G be a Hilbert space over K and let \(v \in \mathcal{P}(\mathrm{F};\mathrm{G})\) be an operator with dense domain. If \(v \circ u\) has dense domain, then \(u^{*} \circ v^{*} \subset (v \circ u)^{*}\) . In general, there is no equality (loc. cit.). If u (resp. v) is an isomorphism, then \(v \circ u\) has dense domain and one has \((v \circ u)^{*} = u^{*} \circ v^{*}\) . This is the case, for example, if E = F (resp. F = G) and \(u = \lambda 1_{E}\) (resp. \(v = \lambda 1_{F}\) ) where \(\lambda \in K^{*}\) .

We denote \(s\) the isometric isomorphism of Hilbert spaces from \(\mathrm{E} \oplus \mathrm{F}\) in \(\mathrm{F} \oplus \mathrm{E}\) defined by \(s(x, y) = (-y, x)\) for all \((x, y) \in \mathrm{E} \oplus \mathrm{F}\) .

PROPOSITION 7. --- Let u be a partial operator with dense domain from E into F.

\begin{enumerate}
\def\labelenumi{\alph{enumi})}
\item
  The graph of \(u^{*}\) is equal to \(s(\Gamma_{u}^{\circ}) = s(\Gamma_{u})^{\circ}\) ;
\item
  The partial operator \(u^{*}\) is closed;
\item
  The kernel of \(u^{*}\) is the orthogonal complement of the image of \(u\) .
\end{enumerate}

Denote by \(\mathrm{W} = s(\Gamma_u)^\circ\) . Since the linear map \(s\) is unitary, we have \(\mathrm{W} = s(\Gamma_u^\circ)\) .

One has \((y,x)\in \mathrm{W}\) if and only if

\[
\langle (y, x) \mid (- u (x ^ {\prime}), x ^ {\prime}) \rangle = 0
\]

for all \(x' \in \mathrm{dom}(u)\) , that is, if

\[
\langle y \mid u (x ^ {\prime}) \rangle = \langle x \mid x ^ {\prime} \rangle
\]

for all \(x' \in \text{dom}(u)\) When \(y \in \text{dom}(u^{*})\) and \(x = u^{*}(y)\) , this property is true (cf.formula (1), p. 236). Conversely, if this condition holds, we deduce that \(|\langle y \mid u(x') \rangle| \leqslant \|x\| \|x'\|\) for all \(x' \in \text{dom}(u)\) , which implies that \(y\) belongs to \(\text{dom}(u^*)\) ; we then have \(u^*(y) = x\) . Hence \(W = \Gamma_{u^*}\) .

The operator \(u^{*}\) is closed, because the space \(s(\Gamma_{u})^{\circ}\) is closed in \(\mathrm{F} \oplus \mathrm{E}\) .

Let us prove assertion \(c\) ). If \(y\) is orthogonal to the image of \(u\) , the linear form \(\lambda_y: x \mapsto \langle y | u(x) \rangle\) on D is zero, hence \(y \in \text{dom}(u^*)\) and \(u^*(y) = 0\) . Conversely, let \(y \in \text{dom}(u^*)\) . We then have \(u^*(y) = 0\) if and only if \(y\) is orthogonal to \(u(x)\) for all \(x \in D\) (formula (1), p. 236).

PROPOSITION 8. --- Let u be an operator with dense domain from E to F. Then \(u^{*}\) has dense domain if and only if u is closable. When this is the case, the closure \(\overline{u}\) of u is equal to \(u^{**}\) , and the adjoint of \(\overline{u}\) is equal to \(u^{*}\) .

By prop. 7, the partial operator \(u^*\) is closed. Suppose that the domain D* of \(u^*\) is dense in F. Let \(u^{**}\) be the adjoint of \(u^*\) ; it is a closed partial operator from E to F. Let us prove that \(u \subset u^{**}\) , which will imply that \(u\) is closable. Let \(x \in \text{dom}(u)\) By definition of \(u^*\) , the linear forms on D* given by \(y \mapsto \langle x | u^*(y) \rangle\) and \(y \mapsto \langle u(x) | y \rangle\) are equal; we therefore have \(x \in \text{dom}(u^{**})\) and \(u^{**}(x) = u(x)\) whence the assertion.

Conversely, suppose that \(u\) is closable; we have \(\Gamma_{\overline{u}} = \overline{\Gamma}_u\) (prop. 1 of IV, p. 228). Let \(y \in F\) be a vector orthogonal to \(\mathrm{dom}(u^*)\) The element \((y,0)\) of \(F \oplus E\) then belongs to the orthogonal of the graph of \(u^*\) . Now, by prop. 7, a), we have

\[
\Gamma_ {u ^ {*}} ^ {\circ} = (s (\Gamma_ {u}) ^ {\circ}) ^ {\circ} = \overline {{s (\Gamma_ {u})}} = s (\overline {{\Gamma}} _ {u})
\]

It follows therefore that \((0,y)\in\Gamma_{\overline{u}}\) , whence \(y=\overline{u}(0)=0\) . The orthogonal of dom \((u^{*})\) being reduced to 0, the space dom \((u^{*})\) is dense in F.

Finally, prop. 7, applied to \(u^{*}\) implies that

\[
\Gamma_ {u ^ {* *}} = s ^ {- 1} (\Gamma_ {u ^ {*}} ^ {\circ}) = s ^ {- 1} (s (\Gamma_ {u} ^ {\circ \circ})) = \overline {{\Gamma}} _ {u},
\]

hence \(u^{**} = \overline{u}\) , then \(\overline{u}^{*} = (u^{*})^{**} = \overline{u^{*}} = u^{*}\) as \(u^{*}\) is closed.

COROLLARY. --- If u is a closed partial operator with dense domain from E into F, then \(u^{*}\) has dense domain and one has \(u^{**} = u\) .

DEFINITION 7. --- Let u be a partial operator on E. One says that u is symmetric if u is densely defined and if \(u^{*}\) is an extension of u. One says that u is self-adjoint if u is densely defined and \(u^{*} = u\) .

u is said to be essentially self-adjoint if it is closable and if its closure \(\overline{u}\) is self-adjoint.

We say that u is a partial operator bounded below operator if u is symmetric and if there exists a real number c such that \(\langle x | u(x) \rangle \geqslant c \|x\|^{2}\) for all x belonging to the domain of u. One then says that c is a lower bound of u. If c = 0, one also says that u is a positive partial operator.

We denote \(\mathcal{A}(\mathrm{E})\) the set of partial self-adjoint operators on E.

Remarks. --- 1) For an operator \(u\) to be densely defined on E and symmetric, it is necessary and sufficient that one have

\[
\langle x \mid u (y) \rangle = \langle u (x) \mid y \rangle\tag{2}
\]

for all \((x,y)\in\mathrm{dom}(u)^{2}\) This formula indeed shows that the domain of u is contained in that of \(u^{*}\) then that \(u^{*}\) and u coincide on the domain of u. In particular, it follows that \(\langle x|u(x)\rangle\in\mathbf{R}\) for all \(x\in\mathrm{dom}(u)\) .

As will be seen in various examples, formula (2) can often be verified by a straightforward calculation. On the other hand, the exact determination of the domain of the adjoint, which alone makes it possible to know whether a symmetric operator is self-adjoint or not, can be very delicate.

\begin{enumerate}
\def\labelenumi{\arabic{enumi})}
\setcounter{enumi}{1}
\item
  A self-adjoint partial operator u is essentially self-adjoint (cf.prop. 8).
\item
  Let \(u\) a symmetric partial operator on E. The operator \(u\) is closable (prop. 7, b)). It satisfies \(\text{dom}(u) \subset \text{dom}(u^*)\) , and \(u\) is self-adjoint if and only if \(\text{dom}(u) = \text{dom}(u^*)\) . Moreover, the closure \(\overline{u}\) of \(u\) is symmetric since \(\overline{u} \subset u^* = \overline{u}^*\) (prop. 8).
\item
  Suppose K = C. Let \(u \in \mathcal{P}(\mathrm{E};\mathrm{E})\) a densely defined partial operator. The condition \(\langle x | u(x) \rangle \in \mathbf{R}\) for all \(x \in \operatorname{dom}(u)\) implies that u is symmetric (EVT, V, p. 2, remark); in particular, if \(\langle x | u(x) \rangle \in \mathbf{R}_{+}\) for all \(x \in \operatorname{dom}(u)\) , then u is positive.
\item
  Let u and v be symmetric partial operators on E. If u is self-adjoint and if \(u \subset v\) , then \(v \subset v^{*} \subset u^{*} = u\) therefore u = v.
\item
  An essentially self-adjoint partial operator \(u\) is symmetric, since \(u \subset \overline{u}\) implies \(\overline{u} = \overline{u}^{*} \subset u^{*}\) , hence \(u \subset u^{*}\) .
\item
  Let u and v be symmetric partial operators on E. If \(u+v\) is has dense domain, for example if u or v belongs to \(\mathcal{L}(\mathrm{E})\) , then \(u+v\) is symmetric. In general, the partial operator \(u + v\) is not self-adjoint, even if u and v are (exercise 9 of IV, p. 347).
\item
  Let \(u\) a symmetric partial operator on E. A real number \(c\) is a lower bound of \(u\) if and only if the operator \(u - c \cdot 1_{\mathrm{E}}\) is positive.
\end{enumerate}

Lemma 3. --- Suppose that K = R. Let u be a partial operator with dense domain from E into F.

\begin{enumerate}
\def\labelenumi{\alph{enumi})}
\item
The adjoint of \(u_{(\mathbf C)}\) is \((u^*)_{(\mathbf C)}\) ;
\item
  Suppose that E = F; the partial operator u is symmetric (resp. self-adjoint) if and only if the partial operator \(u_{(\mathbf{C})}\) is symmetric (resp. self-adjoint).
\end{enumerate}

Let us prove \(a\) ). Let \(y \in \mathrm{F}_{(\mathbf{C})}\) and let us write \(y = y_1 + iy_2\) with \(y_1, y_2 \in \mathrm{F}\) . For all \((x_1, x_2) \in \mathrm{E} \times \mathrm{E}\) , we have

\[
\begin{array}{c} \langle u _ {(\mathbf {C})} (x _ {1} + i x _ {2})   |   y \rangle = \langle u (x _ {1})   |   y _ {1} \rangle + i \langle u (x _ {1})   |   y _ {2} \rangle \\ - i \langle u (x _ {2})   |   y _ {1} \rangle + \langle u (x _ {2})   |   y _ {2} \rangle . \end{array}
\]

If \(y \in \text{dom}(u^*)(\mathbf{C})\) , we deduce that \(y \in \text{dom}((u_{(\mathbf{C})})^*)\) and that \(u_{(\mathbf{C})}^*(y) = (u_{(\mathbf{C})})^*(y)\) , hence \(u_{(\mathbf{C})}^* \subset (u_{(\mathbf{C})})^*\) .

Conversely, suppose that \(y\) belongs to \(\text{dom}((u_{(\mathbf{C})})^*)\) . If one takes \(x_2 = 0\) (resp. \(x_1 = 0\) ) in the above formula, one verifies that \(y_1 \in \text{dom}(u^*)\) (resp. that \(y_2 \in \text{dom}(u^*)\) ), whence \(y \in \text{dom}(u^*)(\mathbf{C})\) .

The assertion \(a\) ) implies that \(u_{(\mathbf{C})}\) is symmetric (resp. self-adjoint) if \(u\) is.

Conversely, suppose that \(u_{(\mathbf{C})}\) is symmetric. The relation \(\langle u(x)|y\rangle = \langle x|u(y)\rangle\) for all \((x,y) \in \mathrm{dom}(u_{(\mathbf{C})}) \times \mathrm{dom}(u_{(\mathbf{C})})\) implies that u is symmetric by taking x and y in the subspace \(\mathrm{dom}(u)\) of \(\mathrm{dom}(u_{(\mathbf{C})})\) . If \(u_{(\mathbf{C})}\) is self-adjoint, what precedes shows that u is symmetric; since \(\mathrm{dom}(u^{*}) = \mathrm{dom}(u_{(\mathbf{C})}^{*}) \cap \mathrm{F}\) , assertion a) implies that \(\mathrm{dom}(u^{*}) = \mathrm{dom}(u_{(\mathbf{C})}) \cap \mathrm{F} = \mathrm{dom}(u)\) , hence u is self-adjoint.

